\section{Contexto del problema}

Un hospital público dispone de información obtenida durante controles cardiovasculares preventivos, pero no cuenta con una solución analítica que transforme sistemáticamente esos datos en conocimiento sobre los factores asociados a enfermedad cardíaca ni que apoye la priorización de acciones preventivas.

La solución propuesta es una plataforma interna de apoyo analítico para prevención cardiovascular. Debe procesar información clínica, analizar factores asociados a enfermedad cardíaca y evolucionar hacia estimaciones de riesgo que ayuden a priorizar evaluaciones preventivas. La supervisión de profesionales de salud se mantiene durante todo el proceso.

\subsection{Objetivo general}

Diseñar y evaluar una solución de ciencia de datos mantenible, segura y explicable para analizar factores asociados a enfermedad cardíaca y apoyar la priorización de pacientes en controles cardiovasculares preventivos. La evaluación compara alternativas Cloud, Local e Híbrida, sin delegar la decisión clínica en el sistema.

\subsection{Objetivos específicos}

\begin{itemize}
  \item Preparar y analizar los datos cardiovasculares para determinar su calidad, representatividad y utilidad analítica.
  \item Identificar variables y patrones asociados con la presencia de enfermedad cardíaca.
  \item Desarrollar y evaluar progresivamente modelos de apoyo a la estimación de riesgo cardiovascular.
  \item Incorporar explicabilidad, evaluación de sesgos, privacidad y supervisión humana en el ciclo de la solución.
  \item Diseñar una arquitectura mantenible que permita incorporar nuevos datos, monitorear la solución y evolucionarla.
\end{itemize}

\subsection{Usuarios y criterios de éxito}

Los profesionales clínicos y preventivos consultarán resultados como apoyo; coordinación y gestión cardiovascular utilizarán información agregada para planificar acciones; y los equipos de TI y datos administrarán operación, permisos, monitoreo y soporte bajo el principio de mínimo privilegio. El paciente es un beneficiario indirecto y no se considera una interfaz pública de paciente para el piloto.

El éxito técnico requiere una arquitectura mantenible y trazable; el éxito clínico-operacional exige resultados comprensibles sin bloquear el flujo clínico si la plataforma no está disponible; el éxito ético requiere explicabilidad, supervisión humana, protección de datos y evaluación por subgrupos; y el éxito financiero y de gestión exige una alternativa sostenible, con alcance, costos, riesgos, recursos y decisiones coherentes.
